\section{Área Problemática}

\subsection{Definición del Problema Central}

El problema central investigado es la incapacidad de los sistemas actuales de gestión
de la calidad para detectar automáticamente, en tiempo real, picos emocionales en las
llamadas de un call center. El proceso estándar exige que un supervisor escuche
grabaciones completas: una llamada de 90 minutos requiere 90 minutos de análisis, lo
que limita la cobertura real al 1--5\,\% del total de interacciones \cite{zoom2025}.

\subsection{Contexto e Impacto del Problema}

El sector de call centers está valuado en más de 400 mil millones de dólares, con una
proyección de 496 mil millones para 2027 \cite{zoom2025}. En Colombia y en América
Latina, los call centers constituyen uno de los sectores de mayor empleo del sector
de servicios. El uso de \textit{speech analytics} creció del 28\,\% al 37{,}5\,\%
entre 2022 y 2023 \cite{zoom2025}, lo que evidencia una demanda creciente de
automatización que los sistemas actuales no logran satisfacer por completo.

El impacto se manifiesta en tres dimensiones: (a) \textbf{Operacional:} solo el 31\,\%
de los contact centers mide activamente las emociones del cliente \cite{zoom2025}.
(b) \textbf{Humana:} el 88\,\% de los profesionales del sector reconoce el burnout de
los agentes como uno de los mayores desafíos de la industria, con una tasa promedio de
rotación anual entre el 38\,\% y el 52\,\% \cite{sqm2023, hiringbranch2024}.
(c) \textbf{Calidad de servicio:} el número de consumidores que abandonan una marca
tras una sola interacción negativa aumentó en 10 puntos porcentuales entre 2022 y
2024 \cite{zoom2025}. La Tabla~\ref{tab:impacto} resume los indicadores clave.

\begin{table}[!t]
\renewcommand{\arraystretch}{1.3}
\caption{Indicadores de Impacto del Problema}
\label{tab:impacto}
\centering
\begin{tabular}{l p{3.5cm} p{2cm}}
\toprule
\textbf{Dimensión} & \textbf{Indicador} & \textbf{Dato} \\
\midrule
Operacional & Centros que miden emociones activamente                   & Solo el 31\,\% \cite{zoom2025} \\
Operacional & Crecimiento uso de \textit{speech analytics} 2022--2023  & 28\,\% $\to$ 37{,}5\,\% \cite{zoom2025} \\
Humana      & Agentes con alto nivel de burnout (SQM, 2022)            & 63\,\% \cite{sqm2023} \\
Humana      & Tasa de rotación anual promedio (Deloitte, 2023)         & 38--52\,\% \cite{hiringbranch2024} \\
Humana      & Agentes que reportan tensión laboral alta (Cornell)      & 87\,\% \cite{hiringbranch2024} \\
Calidad     & Clientes que abandonan la marca tras 1 interacción negativa & +10 pp (2022--2024) \cite{zoom2025} \\
\bottomrule
\end{tabular}
\end{table}

\subsection{Preguntas Clave de Investigación}

\textbf{¿Cuál es el problema que se investigará?} Se investigará cómo detectar
automáticamente picos emocionales en llamadas de call center mediante el análisis de
características paralingüísticas de la voz —pitch, intensidad, jitter, shimmer y
MFCCs— y algoritmos de machine learning, sin depender del contenido textual de la
conversación.

\textbf{¿A quiénes afecta y en qué medida?} El problema afecta a tres grupos. Los
agentes de call center son el grupo más directamente perjudicado: el 76\,\% de los
practicantes del sector reconoce altos niveles de burnout \cite{sqm2023} y el 87\,\%
de los agentes reporta una alta o muy alta tensión laboral \cite{hiringbranch2024}.
Los supervisores de calidad están limitados a revisar entre el 1\,\% y el 5\,\% de
las llamadas. Las organizaciones asumen costos de rotación del 30\,\% al 50\,\% del
salario anual de un agente por cada reemplazo \cite{sqm2023}.

\textbf{¿Cuáles son las causas que contribuyen al problema?} Las causas se agrupan en
tres niveles. (a) \textbf{Tecnológico:} los sistemas de análisis de sentimiento
basados en texto pierden información prosódica y paralingüística esencial
\cite{feng2023, jha2022}. (b) \textbf{Metodológico:} la ausencia de baselines
emocionales individualizadas hace que los sistemas con umbrales fijos produzcan tasas
elevadas de falsos positivos y falsos negativos \cite{gat2022}. (c)
\textbf{Organizacional:} la cobertura de monitoreo manual es estructuralmente
insuficiente, pues el tiempo disponible para la revisión es igual al de cada llamada.

\textbf{¿Qué consecuencias tiene no intervenir?} La no intervención perpetúa un ciclo
de degradación simultáneo: para los agentes, la exposición crónica sin apoyo oportuno
acelera el burnout y la rotación —en 2023 alcanzó el 52\,\% según Deloitte
\cite{hiringbranch2024}—. Para las organizaciones, cada punto porcentual de rotación
no mitigado implica costos directos de contratación y pérdida de conocimiento
institucional. Para los clientes, la falta de detección impide una intervención
supervisora oportuna, lo que deteriora la experiencia del usuario \cite{zoom2025}.
